%!TEX program = xelatex
\documentclass{standalone}
\usepackage{ifpdf}
\usepackage{tikz}
\usepackage[active,tightpage]{preview}
\PreviewEnvironment{tikzpicture}

\usepackage{tikz}
\usetikzlibrary{arrows.meta, backgrounds, fit, positioning, shapes.geometric, calc, decorations.pathreplacing, decorations.markings,shadows.blur}

\usepackage{xltxtra}
\usepackage{xunicode}

% Fonts
\usepackage{fontspec}
\defaultfontfeatures{Mapping=tex-text}

\usepackage{newtxsf}
\setmainfont{Fira Sans}
\usepackage{mathastext}

\begin{document}

\begin{tikzpicture}[
        node distance=1.25cm,
        box/.style={rectangle, draw, thick, minimum width=2cm, minimum height=1cm, align=center, rounded corners=3pt},
        data/.style={box, fill=blue!20},
        algorithm/.style={box, fill=red!20},
        output/.style={box, fill=green!20},
        arrow/.style={-{Stealth[scale=1.2]}, thick},
        title/.style={font=\bfseries},
        subtitle/.style={font=\footnotesize\bfseries}
    ]

    % Classical Programming Paradigm
    \node[title] (title1) {\underline{Classical Programming Paradigm}};
    \node[data, below=0.8cm of title1] (data1) {Data};
    \node[algorithm, right=of data1] (algorithm1) {Algorithm\\(Rules)};
    \node[output, right=of algorithm1] (output1) {Output\\(Results)};

    \draw[arrow] (data1) -- (algorithm1);
    \draw[arrow] (algorithm1) -- (output1);

    \node [below=0.3cm of data1, subtitle] {Known};
    \node [below=0.3cm of algorithm1, subtitle] {\sc Designed by programmer};
    \node [below=0.3cm of output1, subtitle] {Desired};

    \begin{scope}[on background layer]
        \node[fit=(title1)(data1)(algorithm1)(output1), rectangle, draw=gray!30, rounded corners, inner sep=0.8cm, fill=gray!5] {};
    \end{scope}

\end{tikzpicture}

\end{document}
